\subsection{Background}

Directional antennas are frequently used in wireless communications to enhance the distance, SNR, and reduce interference by focusing the radiation power towards a certain direction. Directional antennas are often utilized in fixed wireless connections, point-to-point backhauls, wireless sensor networks, and long-distance IoT networks. However, the effectiveness of directional antennas depends heavily on precise alignment between the transmitting antenna and the receiving antenna.

However, the problem of optimal antenna orientation involves several key issues, among which are the effects of multipath fading, the movement of environmental objects, vibration, and changes in antenna mount position. Antenna misorientation results in severe degradation of signal strength and, ultimately, system efficiency. Manual adjustment of antennas has always been required in the past, but it is often inconvenient and unfeasible in applications where continual operation is necessary. Existing automatic techniques typically employ expensive hardware such as phase array antennas, angle of arrival estimation algorithms, and complex RF frontend signal processing, which results in increased cost, power consumption, and complexity of design. Such solutions cannot be utilized in inexpensive embedded systems. Thus, the development of receiver-based alignment approaches becomes increasingly popular.

The current approaches to automating the process of alignment usually rely on high-cost hardware technologies, such as phased antenna arrays, angle-of-arrival determination, or advanced RF front end processing techniques. This increases the cost and design complexity of the systems, making their use in embedded, low-cost applications unfeasible. As a result, more attention is being devoted to developing autonomous alignment methods at the receiver side using easily measurable signal quality parameters, such as RSSI and throughput.